\section{Results and Discussion}
\label{sec:results}

\subsection{Campaign Validity}

The main campaign executed \NumRunsMain{} runs over \NumContexts{} contexts with
\NumPolicies{} policies and \NumSeeds{} paired seeds. No run failed. Every run
completed every vehicle it generated (minimum completion rate
\NumMinCompletion{}), with \NumTeleports{} teleported vehicles in the entire
campaign, so no result rests on a vehicle the simulator removed. All
\NumContextsAnalysed{} contexts satisfy the validity and completeness rules
declared in advance, and none is excluded.

\subsection{Policy Complementarity}
\label{sec:complementarity}

The confirmatory criterion G1$'$ \NumGOneVerdict{}: \NumGOneNPolicies{} policies
(\NumGOnePolicies) are strictly best in at least one context with a
seed-resolved margin. \textbf{P2 is never a resolved winner anywhere in the
factor space.} Resolved winner counts are P3 \NumWinResPThree{}, P4
\NumWinResPFour{} and P1 \NumWinResPOne{} (Table~\ref{tab:complementarity}).

The winner map is multi-factor but shallow. Naming the same policy everywhere is
correct in \NumWinMapConst{} of contexts; the best single factor
(\NumWinMapOneFactor) reaches \NumWinMapOne{}, and the best pair of factors
\NumWinMapTwo{}. The structure is interpretable and matches the design
(Figure~\ref{fig:regime}): P1 wins only where the shortest corridor is
uncongested \emph{and} penetration is high enough for the adaptive policies to
over-divert; P4 wins at low demand with generous green; P3 wins every context
at the highest demand level.

This answers RQ1 affirmatively but narrowly. The portfolio is complementary in
the sense the algorithm-selection literature uses --- more than one policy is
uniquely best somewhere, and the map depends on more than one factor --- yet one
policy accounts for the large majority of wins. That combination is what makes
the next two subsections necessary.

\subsection{Decision Value of Adaptation}

Complementarity and decision value are different quantities, and here they point
in opposite directions.

The single-best fixed policy is \NumGlobalSBS{}, optimal in \NumSBSOptimalPct{}
of contexts. Figure~\ref{fig:performance} shows why the portfolio looks the way
it does: the distance-minimising policy diverges with demand while the three
adaptive policies track one another closely. Running the single-best policy
always costs \NumSBSMeanCost~s of system mean journey time per vehicle; the
retrospective per-context best costs \NumCostVBS~s. The headroom of
Eq.~\ref{eq:headroom} is therefore
\begin{equation}
H \;=\; \NumHeadroom\ \text{s per vehicle},
\qquad
\frac{H}{\mathbb{E}_x[C(x,a_{\mathrm{SBS}})]}
\;=\; \frac{\NumHeadroom}{\NumSBSMeanCost} \;=\; \NumHeadroomPct .
\end{equation}
This is the aggregate decision headroom relative to the single-best fixed
policy. It is not a claim that any vehicle's travel time improves by
\NumHeadroomPct{}; it is the share of the fixed policy's achieved cost that a
perfect per-context choice could remove.

Selecting the wrong \emph{fixed} policy costs far more. Using the
distance-minimising policy throughout increases system mean journey time by
\NumFixedPenAbsPOne~s per vehicle relative to the single-best fixed policy,
which is \NumFixedPenPctAbsPOne{} of that policy's mean journey time of
\NumSBSMeanCost~s. The reactive and reliability-aware policies cost
\NumFixedPenAbsPTwo~s (\NumFixedPenPctAbsPTwo{}) and \NumFixedPenAbsPFour~s
(\NumFixedPenPctAbsPFour{}) on the same basis.

\begin{quote}
Choosing the fixed policy well is worth \NumFixedPenPctAbsPOne{} of mean journey
time. Choosing it \emph{per context}, given that the best fixed policy is
already in use, is worth \NumHeadroomPct{}.
\end{quote}

The two decisions differ by roughly two and a half orders of magnitude. A
benchmark that reports only the second makes the first invisible, and a system
built to exploit the second while getting the first wrong would be worse than
no system at all.

\subsection{Why the Adaptation Headroom Is Small}
\label{sec:asymmetry}

The explanation is not that the policies behave alike; Section~\ref{sec:%
complementarity} shows they do not. It is that the best fixed policy is
near-dominant in a measurable sense.

Define the \emph{downside} of the fixed policy as its mean loss in those
contexts where it is not best, and its \emph{upside} as its mean gain over the
best alternative in those contexts where it is. In \NumAsymNotBestPct{} of
contexts \NumGlobalSBS{} is not best, and there it loses \NumAsymLoss~s on
average, at most \NumAsymMaxLoss~s. In \NumAsymBestPct{} of contexts it is best,
and there it wins by \NumAsymMargin~s on average, at most \NumAsymMaxMargin~s.
Its upside is \textbf{\NumAsymRatio$\times$ its downside}.

A portfolio containing such a member has a large ranking structure and a small
decision structure: any selector moving away from it risks a large loss to
chase a small gain. This is an empirical characterisation of this portfolio in
this environment, not a general theorem. It is, however, measurable in any
portfolio before a selector is built, and we suggest measuring it first.

\subsection{Noise and Statistical Resolution}

Among the \NumResolvNContexts{} contexts where the fixed policy is not best, the
mean headroom is \NumResolvHeadroomNZ~s against a mean two-standard-error band
of \NumResolvNoiseNZ~s on the same paired comparison, and the headroom exceeds
its own noise band in \NumResolvExceeds{} of them (\NumResolvExceedsPct{}).
Resolving the campaign-wide mean effect of \NumResolvHeadroom~s under the
criterion of Eq.~\ref{eq:resolved} would require roughly
\NumResolvSeedsNeeded{} seeds per context, about \NumResolvRunsNeeded{} runs.

We report this rather than presenting a near-noise-floor improvement as an
achievement. It also bounds what the remaining subsections can claim: selector
differences of a few hundredths of a second per vehicle are real as point
estimates and are not separable from replication noise context by context.

\subsection{Mechanistic Rule}

Fitted on all contexts for exposition, the depth-2 rule splits first on
$\NumMechFeature = p\,D/\mathrm{cap}_C$ --- the share of the shortest corridor's
capacity that the guided cohort alone would consume --- at a threshold bracketed
between \NumMechBracketLo{} and \NumMechBracketHi{}. The rule is physically
readable: below that share the decision turns on the effective green ratio;
above it, on whether the arterial is already oversaturated
(Figure~\ref{fig:regions}). This answers RQ2: the preference boundary is
explicable in mechanistic terms, and the explanatory variable is a degree of
saturation rather than a raw demand.

The rule identifies the best policy in \NumMechTopOne{} of contexts, more often
than always naming the fixed policy does. Its mean decision regret is
nevertheless \NumMechRegret~s against \NumMechFixedRegret~s for the fixed
policy. \textbf{It is more accurate and more costly.} Under
leave-one-context-out the same pattern holds: B1 reaches \NumBOneTopOne{} top-1
accuracy against B0's \NumBZeroTopOne{}, and \NumBOneRegret~s of regret against
\NumBZeroRegret~s.

This is the predict-then-optimise distinction~\cite{elmachtoub2022spo} in a
traffic setting. A rule trained to name the winner is rewarded for being right
where the margin is negligible and punished only mildly for being wrong where it
is large --- the wrong trade under an \NumAsymRatio$\times$ asymmetry. The
practical implication is that selection accuracy should not be reported as a
proxy for decision quality; in this study the two criteria rank the methods
differently.

\subsection{Advantage-Based Selector}

Predicting advantage (Eq.~\ref{eq:advantage}) rather than identity reverses the
failure. B4 reaches \NumBFourRegret~s of mean regret, closing
\NumBFourGapClosed{} of the gap between the fixed policy and the retrospective
best, while being \emph{less} accurate at naming the winner (\NumBFourTopOne{})
than the logistic baseline (\NumBThreeTopOne{}, \NumBThreeRegret~s). It also has
the lowest maximum regret of any selector, \NumBFourMax~s against
\NumBZeroMax~s (Table~\ref{tab:ladder}, Figure~\ref{fig:ladder}).

The primary comparison declared in advance was B4 against B1. B4 has
substantially lower mean decision regret than B1 (\NumBFourRegret{} versus
\NumBOneRegret~s per vehicle). Beating B0 alone would not have been evidence of
much, since B0 is near-dominant by Section~\ref{sec:asymmetry}. This answers
RQ3: a learned selector does identify structure beyond the mechanistic rule ---
but only when trained on the decision cost rather than on the winner's identity.

The magnitude is what matters operationally. B4 reduces mean regret by
\NumBFourVsBZero~s per vehicle relative to the fixed policy, which is
\NumBFourVsBZeroPct{} of a \NumSBSMeanCost~s mean journey time. The selector is
real and reproducible, and below the threshold at which an operator would act on
it.

\subsection{Selective Prediction and Distribution Shift}

B5 abstains in \NumAbstainRate{} of in-distribution contexts and returns exactly
the fixed policy's regret, \NumBFiveRegret~s. The mechanism is explicit rather
than tuned: the conformal interval on the predicted advantage is wider than the
advantage itself, so the requirement that its lower bound exceed zero is never
met. A calibrated gate refuses to act on an effect smaller than its own
uncertainty, which is the correct behaviour given Section~6.5.

Across the \NumOODNTotal{} shift splits (Table~\ref{tab:ood},
Figure~\ref{fig:ood}) the unguarded selector B4 is \textbf{not uniformly
reliable}. Using a tolerance of \NumOODTolerance~s per vehicle, it helps on
\NumOODNHelps{} splits, is indistinguishable on \NumOODNNeutral{}, and harms on
\NumOODNHarms{}:
\begin{itemize}
\item \emph{helps}: the interpolation control O3 (\NumOODThreeMidBZero{} to
\NumOODThreeMidBFour~s) and the domain shift O8 (\NumOODEightNorthBZero{} to
\NumOODEightNorthBFour~s);
\item \emph{indistinguishable}: unseen low demand O2, unseen information lag
O4, unseen penetration O5;
\item \emph{harms}: unseen high demand O1 (\NumOODOneHighBZero{} to
\NumOODOneHighBFour~s), unseen green ratio O6, and most severely the concept
shift O7 (\NumOODSevenIncBZero{} to \NumOODSevenIncBFour~s).
\end{itemize}
Harmful cases therefore occurred under both extrapolation and concept shift,
while two splits improved. \textbf{On each of the \NumOODNHarms{} harmful splits,
B5 recovered the fixed policy's regret exactly.} The selective rule prevented
the harms that were observed here; it does not guarantee protection against
harms not observed here.

The support gate flags \NumOODFourLagSupport{}, \NumOODFivePenSupport{} and
\NumOODSixGreenSupport{} of held-out contexts as outside the training envelope
on the lag, penetration and green-ratio extrapolations respectively. The
mechanistic rule degrades worst of all under shift, reaching
\NumOODFivePenBOne~s on unseen penetration, which is consistent with a rule
whose thresholds were fitted inside a range it is now being asked to leave.

\paragraph{A shift the gate cannot see.} O8 is instructive precisely because the
support gate does \emph{not} fire: it flags \NumOODEightNorthSupport{} of
held-out contexts, because those contexts' descriptors are identical to
descriptors seen in training --- the incident's \emph{location} is not among the
eight. The selector happened to do well there, but it did so unguarded. A
support gate defined on a feature space cannot detect a shift that leaves that
feature space unchanged, and this is a concrete instance rather than a
theoretical caveat. It answers RQ4 with a qualification: abstention removed
every harm we were able to observe, and its coverage is bounded by what the
feature representation can express.

\subsection{Boundary Refinement}
\label{sec:boundary}

\NumBoundMono{} of the \NumBoundN{} refined sequences are monotone, with a
single change point, and in point estimate they narrow the switching bracket
from the \NumBoundGrid~veh/h grid spacing to a median of \NumBoundWidth~veh/h.
All \NumBoundPOneN{} shortest-to-load-balancing transitions are monotone; their
change points sit in a \NumBoundWidth~veh/h bracket whose position differs by
cell, from \NumBoundPOneEarliestLo--\NumBoundPOneEarliestHi~veh/h to
\NumBoundPOneLatestLo--\NumBoundPOneLatestHi~veh/h. That the position moves with
the cell is the point: the boundary is a property of the operating conditions,
not a single number for the network.

\textbf{None of these refined steps is seed-resolved.} The point estimates move
consistently, but no individual step between adjacent refined levels satisfies
Eq.~\ref{eq:resolved} at \NumSeeds{} seeds. The remaining \NumBoundNonMono{}
sequences are non-monotone, and all involve the reliability-aware policy. We
therefore report the transition as a point-estimate bracket and state that it is
unresolved: the shortest-path to load-balancing reversal is an ordered regime
change, while the apparent reliability-aware transitions near it are not
distinguishable from noise.

\subsection{Multi-Criteria and Pareto Structure}
\label{sec:pareto}

Over all \NumContextsAnalysed{} contexts, \NumParetoNonSingleton{} have a
non-singleton Pareto set. Non-domination is not concentrated in one policy: P3
is non-dominated in \NumParetoPctPThree{} of contexts, P4 in
\NumParetoPctPFour{}, P2 in \NumParetoPctPTwo{} and P1 in \NumParetoPctPOne{}.

The criteria are not interchangeable. Journey time and \CO{} correlate at
\NumCorrCOneCTwo{} and rank the four policies identically in
\NumSameOrderCOneCTwo{} of contexts; journey time and stopped delay correlate at
\NumCorrCOneCThree{} and agree in only \NumSameOrderCOneCThree{}. The best
policy on \CO{} differs from the best on journey time in
\NumBestCTwoDiffersPct{} of contexts; on stopped delay it differs in
\NumBestCThreeDiffersPct{}.

The policies specialise by criterion rather than by regime.
\NumSpecCOnePolicy{} minimises journey time in \NumSpecCOneN{} contexts and
\NumSpecCTwoPolicy{} minimises \CO{} in \NumSpecCTwoN{}, but
\textbf{\NumSpecCThreePolicy{} minimises stopped delay in \NumSpecCThreeN{}
contexts (\NumSpecCThreePct{})} while minimising journey time in only
\NumSpecCOnePFour{}. The reliability-aware policy routes towards the steady,
unsignalised corridor: its vehicles spend less time stationary and more time
moving.

\textbf{This is where the portfolio's complementarity lives} --- on the criterion
vector, not on the scalar cost. An operator whose objective is queue clearance
at signalised junctions and one whose objective is corridor journey time should
not deploy the same policy, and that is a different argmin in
\NumBestCThreeDiffersPct{} of contexts rather than a preference weighting.

\subsection{Negative Controls and Sensitivity Findings}

\paragraph{Eco-routing is excluded on evidence.} \CO{} and journey time
correlate at \NumCorrCOneCTwo{} and select the same policy in
\NumSameOrderCOneCTwo{} of contexts. An objective nearly monotone in travel time
adds a label rather than a policy. The numbers are given so a reader may
disagree with the threshold rather than with an assertion.

\paragraph{A policy that never wins.} P2, reactive dynamic travel-time routing,
is never a resolved winner in any of the \NumContextsAnalysed{} contexts, and is
non-dominated in only \NumParetoPctPTwo{}. In an environment with three
alternatives and an information lag, reacting to the mean is dominated by
balancing load. This is a negative result about a widely deployed policy family
in this environment, and it is reported rather than omitted.

\paragraph{Preference weighting.} Across the three declared preference profiles
the preferred policy changes in \NumPrefChanges{} of contexts
(Table~\ref{tab:sensitivity}). The journey-time choice is reproduced by the
time-weighted profile in \NumPrefAgreeTime{} of contexts, the balanced profile
in \NumPrefAgreeBalanced{} and the environment-weighted profile in
\NumPrefAgreeEnvironment{}. The decision is stable under reasonable preference
assumptions, and where it is not, stopped delay is what moves it.

\paragraph{Policy parameters.} The dispersion weight of the reliability-aware
policy is close to immaterial, changing that policy's own mean journey time by
at most \NumLambdaMaxChangePct{}; across all four parameter changes the identity
of the best policy is unchanged in \NumWinnerStableMin{} to
\NumWinnerStableMax{} of the contexts tested. The admissibility tolerance of the
load-balancing policy is not immaterial: loosening it from \NumEpsFrozen{} to
\NumEpsAlt{} improves that policy's own mean journey time by \NumEpsGainPct{}
(\NumEpsGainSec~s). The value fixed in advance therefore \emph{understates} how
dominant the load-balancing policy is, which strengthens rather than weakens the
central conclusion. We report this rather than adopting the better value
retrospectively.

\subsection{Interpretation, Scope, and Limitations}

\paragraph{What the benchmark shows.} Three questions that are usually treated
as one come apart here. The portfolio is complementary; the decision value of
exploiting that complementarity is \NumHeadroomPct{} of the fixed policy's mean
journey time; and a selector capturing that value is not uniformly reliable
outside its training range. The decision carrying substantial value in this
environment is which fixed policy to run --- worth \NumFixedPenPctAbsPOne{} ---
together with which criterion to optimise, since the policy minimising stopped
delay is not the policy minimising journey time in \NumBestCThreeDiffersPct{} of
contexts. Both the headroom and its noise floor are computable from paired
replications before any model is built, which we suggest as a standing first
step.

\paragraph{Environment and scope.} The network is synthetic: a controlled
instrument built to isolate mechanisms, not a model of any city. Corridor demand
is a single origin--destination pair choosing among three corridors; the four
minor streets couple the routing decision to traffic that is not routed, but
many simultaneously-routed pairs, with the interactions between their choices,
are outside what is tested. No claim of real-world validation is made anywhere
in this paper. What is portable is the method rather than the numbers: the
descriptors are dimensionless, so a boundary is stated as a degree of saturation
and a controlled-flow share, both measurable on a real network.

\paragraph{Decision context and compliance.} The selector is given the operating
conditions; a deployed system would have to forecast them and would inherit the
forecast's error on top of the selector's. Guided vehicles comply fully, and
unguided vehicles never respond to conditions at all; real behaviour lies
between these. The reported decision quality is therefore an upper bound on what
the same selector would achieve behind a forecast and partial compliance.

\paragraph{Resolution and statistical reach.} The experiment samples a finite
grid with \NumSeeds{} seeds. Boundaries are reported as brackets between
adjacent sampled levels and never as continuous thresholds, and no refined
switching step is seed-resolved (Section~\ref{sec:boundary}). Where a factor is
associated with a change of preferred policy, that is an association within a
designed factorial, not an identified causal mechanism. We report the paired
resolution criterion of Eq.~\ref{eq:resolved} and make no claim of statistical
significance in the hypothesis-testing sense.

\paragraph{Uncertainty machinery.} Split conformal coverage requires
exchangeability, which covariate shift breaks; the conformal gate is calibrated
in-distribution and heuristic outside it, the support gate detects only
extrapolation that is visible in the feature space, and neither detects concept
shift. Split O8 is a demonstrated instance of a shift the support gate does not
see.

\paragraph{Fixed parameters and the emission model.} The policy parameters were
fixed before any evaluation run and are not optima; a looser load-balancing
tolerance improves that policy by \NumEpsGainPct{}, which strengthens the
central conclusion but means the reported value is not the best available.
\CO{} is computed by the simulator's HBEFA3 implementation for one vehicle
class, without fleet composition, cold starts or gradient, so C2 should be read
as a consistent comparative index across policies in this environment rather
than as an absolute emissions estimate. No monetary valuation is used anywhere,
which keeps the study free of invented constants but also means it cannot
express a trade-off in currency.

\section{Conclusion}

We studied context-dependent selection among four established urban routing
policies --- shortest distance, reactive travel time, capacity-aware load
balancing and reliability-aware routing --- in \NumRunsMain{} microsimulation
runs over \NumContexts{} controlled contexts, with every policy decision taken
from information available before activation and every comparison paired under
common random numbers.

The portfolio is complementary. Three of four policies are strictly best
somewhere by a seed-resolved margin, the winner map depends on more than one
factor, and \NumParetoNonSingleton{} of contexts have a non-singleton Pareto set
over journey time, \CO{} and stopped delay. The complementarity is strongest on
the criterion vector: the reliability-aware policy minimises stopped delay in
\NumSpecCThreePct{} of contexts while rarely minimising journey time.

The decision value of exploiting it is \NumHeadroom~s of system mean journey
time per vehicle, \NumHeadroomPct{} of the single-best fixed policy's
\NumSBSMeanCost~s, because that policy's average gain where it wins is
\NumAsymRatio$\times$ its average loss where it does not. Selecting the wrong
fixed policy costs \NumFixedPenAbsPOne~s per vehicle, \NumFixedPenPctAbsPOne{}
of the same baseline. Within that narrow margin, a selector trained to predict
the winner's identity is more accurate than the fixed policy and more costly
than it, while one trained to predict each policy's advantage has substantially
lower mean decision regret (\NumBFourRegret{} versus \NumBOneRegret~s per
vehicle) and closes \NumBFourGapClosed{} of the available gap.

Regime coverage, not model capacity, governs whether that value survives
deployment. Across \NumOODNTotal{} distribution-shift splits the unguarded
selector helped on \NumOODNHelps{}, was indistinguishable on \NumOODNNeutral{}
and harmed on \NumOODNHarms{}, with the largest harm under an unseen disruption
regime. A support-and-confidence gate recovered the fixed policy exactly on
every harmful split, at the cost of the in-distribution gain, and one split
demonstrated a shift the support gate cannot detect because the shift leaves the
feature representation unchanged.

These findings are scoped to the synthetic benchmark described here and are not
evidence about any real network. Within that scope, the benchmark separates
three questions usually treated as one --- whether policies are complementary,
whether adaptation carries decision value, and whether a selector can be
deployed safely --- and shows that the answers can differ. The most useful next
step is to repeat the measurement where several origin--destination pairs are
routed simultaneously, since the interaction between their choices is the
mechanism most likely to change the size of the headroom.
